% Luc Destombe --- licence CC BY-NC-SA 4.0
% https://creativecommons.org/licenses/by-nc-sa/4.0/deed.fr
% Fichier autonome : compiler avec pdflatex, deux fois (signets, nombre de pages).
\documentclass[11pt,a4paper]{article}
\makeatletter
\newif\iflycee@unecolonne
\lycee@unecolonnetrue
\newif\iflycee@palatino
\lycee@palatinotrue

\RequirePackage[utf8]{inputenc}
\RequirePackage[T1]{fontenc}
\RequirePackage[french]{babel}
\RequirePackage{etoolbox}

\newcommand{\ShorthandsOff}{\@ifundefined{shorthandoff}{}{\shorthandoff{;:!?}}}
\newcommand{\ShorthandsOn}{\@ifundefined{shorthandon}{}{\shorthandon{;:!?}}}
\ShorthandsOff
\AtBeginDocument{\ShorthandsOn}
\AtBeginEnvironment{tikzpicture}{\ShorthandsOff}

\RequirePackage{geometry}
\RequirePackage{fancyhdr}
\RequirePackage{lastpage}
\RequirePackage{multicol}
\RequirePackage{array}
\RequirePackage{enumitem}
\RequirePackage{xcolor}
\RequirePackage{needspace}

\RequirePackage{amsmath,amssymb}
\RequirePackage{mathpazo}
\renewcommand{\ttdefault}{lmtt}
\DeclareMathSizes{10.95}{10.95}{8}{5.5}
\begingroup\catcode`\_=\active \catcode`\^=\active
\gdef_#1{\sb{\scriptscriptstyle #1}}
\gdef^#1{\sp{\scriptscriptstyle\lycee@moinsepais #1}}
\endgroup
\newcommand\lycee@moins{\mathbin{%
  \setbox\z@\hbox{$\m@th\scriptscriptstyle\lycee@moinsorig$}%
  \dimen@=.55\wd\z@ \advance\dimen@-.5pt
  \hbox to.75\wd\z@{\hss$\m@th\scriptscriptstyle\vcenter{\hbox{%
    \tikz\draw[line width=.5pt,line cap=round](0,0)--(\the\dimen@,0);}}$\hss}}}
\begingroup\lccode`\~=`\- \lowercase{\endgroup\let~}\lycee@moins
\newcommand\lycee@moinsepais{\mathcode`\-="8000 }
\AtBeginDocument{\mathchardef\lycee@moinsorig=\mathcode`\-\relax}
\AtBeginDocument{\catcode`\_=12 \mathcode`\_="8000
  \catcode`\^=12 \mathcode`\^="8000 }
\newcommand\lycee@exposants{%
  \fontdimen13\textfont2=.48\fontdimen6\textfont2
  \fontdimen14\textfont2=.48\fontdimen6\textfont2
  \fontdimen15\textfont2=.48\fontdimen6\textfont2
  \fontdimen13\scriptfont2=.48\fontdimen6\scriptfont2
  \fontdimen14\scriptfont2=.48\fontdimen6\scriptfont2
  \fontdimen15\scriptfont2=.48\fontdimen6\scriptfont2}
\every@math@size\expandafter{\the\every@math@size\lycee@exposants}
\newlength{\retraitcalculs}
\setlength{\retraitcalculs}{2em}

\RequirePackage{tikz}
\usetikzlibrary{arrows.meta,calc,babel,shapes.misc}
\RequirePackage[most]{tcolorbox}
\tcbuselibrary{skins,breakable}

\definecolor{coulDef}{HTML}{1F4E79}
\definecolor{coulProp}{HTML}{7030A0}
\definecolor{coulRem}{HTML}{C55A11}
\definecolor{coulEx}{HTML}{2E7D32}
\definecolor{coulExo}{HTML}{B03A2E}
\definecolor{coulPart}{HTML}{1F4E79}
\definecolor{coulMeth}{HTML}{1B6E4F}
\definecolor{coulCourbe}{HTML}{0B5ED7}
\definecolor{coulCourbeB}{HTML}{C00000}
\definecolor{coulCourbeC}{HTML}{2E7D32}
\definecolor{coulGrille}{HTML}{8C8C8C}
\definecolor{coulRep}{HTML}{1B6E4F}
\definecolor{coulBar}{HTML}{2E7D32}

\newcommand{\R}{\mathbb{R}}
\newcommand{\Rs}{\mathbb{R}^{*}}
\renewcommand{\le}{\leqslant}
\renewcommand{\ge}{\geqslant}
\renewcommand{\approx}{\simeq}
\newcommand{\vect}[1]{\overrightarrow{#1}}
\newcommand{\norme}[1]{\left\lVert #1 \right\rVert}
\newcommand{\intervalle}[2]{\left[#1\,;#2\right]}
\RequirePackage{cancel}
\renewcommand{\CancelColor}{\color{coulExo}}
\newcommand{\fois}[1]{\times{\color{coulExo}#1}}
\newcommand{\barre}[1]{\cancel{#1}}

\tikzset{grille/.style={coulGrille,line width=0.4pt}}
\newcommand{\quadrillage}[4]{\draw[grille,step=1] (#1,#2) grid (#3,#4);}
\tikzset{
  croix/.style={cross out,draw,line width=0.6pt,inner sep=0pt,minimum size=4pt},
  croixdroite/.style={croix,rotate=45,line width=0.8pt},
}
\newcommand{\pt}[3]{\path (#1) node[croix]{} node[#2,font=\small,inner sep=2.5pt]{$#3$};}

\newcommand{\lycee@repere}[4]{%
  \draw[grille,step=1] (#1,#2) grid (#3,#4);
  \draw[-{Stealth[length=2mm]},thick] (#1,0)--({#3+0.4},0) node[below left=-1pt]{$x$};
  \draw[-{Stealth[length=2mm]},thick] (0,#2)--(0,{#4+0.4}) node[below left=-1pt]{$y$};
  \node[below left,font=\scriptsize,inner sep=1.5pt] at (0,0) {$0$};}
\newcommand{\lycee@gradx}[1]{\foreach \k/\l in {#1}{%
  \draw (\k,0.07)--(\k,-0.07) node[below,font=\scriptsize,inner sep=1.5pt]{$\l$};}}
\newcommand{\lycee@grady}[1]{\foreach \k/\l in {#1}{%
  \draw (0.07,\k)--(-0.07,\k) node[left,font=\scriptsize,inner sep=1.5pt]{$\l$};}}
\newcommand{\lycee@intff}[2]{\left[#1\,;#2\right]}
\newcommand{\lycee@intfo}[2]{\left[#1\,;#2\right[}
\newcommand{\lycee@intof}[2]{\left]#1\,;#2\right]}
\newcommand{\lycee@intoo}[2]{\left]#1\,;#2\right[}
\newcommand{\lycee@ens}[1]{\left\{#1\right\}}
\AtBeginDocument{%
  \providecommand{\repere}{\lycee@repere}%
  \providecommand{\gradx}{\lycee@gradx}%
  \providecommand{\grady}{\lycee@grady}%
  \providecommand{\Cf}{\mathcal{C}\sb{\scriptscriptstyle f}}%
  \providecommand{\Cg}{\mathcal{C}\sb{\scriptscriptstyle g}}%
  \providecommand{\intff}{\lycee@intff}%
  \providecommand{\intfo}{\lycee@intfo}%
  \providecommand{\intof}{\lycee@intof}%
  \providecommand{\intoo}{\lycee@intoo}%
  \providecommand{\ens}{\lycee@ens}}

\newlist{questions}{enumerate}{3}
\setlist[questions,1]{label=\arabic*\textdegree),leftmargin=*,itemsep=2pt,topsep=3pt}
\setlist[questions,2]{label=\alph*),leftmargin=*,itemsep=1pt,topsep=2pt}
\setlist[questions,3]{label=\roman*),leftmargin=*,itemsep=1pt,topsep=1pt}
\setlist[itemize]{label=\textbullet,leftmargin=*,itemsep=2pt,topsep=3pt}

\newcommand{\besoinplace}[1]{\par\penalty\@M\begingroup
  \dimen@=#1\relax\dimen@ii=\pagegoal\advance\dimen@ii-\pagetotal
  \ifdim\dimen@>\dimen@ii\ifdim\dimen@ii>\z@\vfil\fi\break\fi\endgroup}

\newcommand{\lycee@pied}{\thepage/\pageref{LastPage}}
\newcommand{\entete}[2]{%
  \pagestyle{fancy}\fancyhf{}%
  \fancyhead[L]{\footnotesize\color{coulPart}\textbf{#1}}%
  \fancyhead[R]{\footnotesize\color{coulPart}\textbf{#2}}%
  \fancyfoot[C]{\footnotesize\lycee@pied}%
  \renewcommand{\headrulewidth}{0.4pt}%
  \renewcommand{\headrule}{\hbox to\headwidth{%
    \color{coulPart!40}\leaders\hrule height \headrulewidth\hfill}}}

\RequirePackage{ccicons}
\newcommand{\lycee@licence}{{\scriptsize\color{gray}%
  \copyright\ Luc Destombe --- \ccbyncsa\ CC BY-NC-SA 4.0}}
\newif\iflycee@licence \lycee@licencetrue
\newcommand{\sanslicence}{\lycee@licencefalse}
\AtBeginDocument{\iflycee@licence\AddToHookNext{shipout/foreground}{%
  \put(0,0){\rlap{\kern\dimexpr1in+\hoffset+\oddsidemargin+\textwidth\relax
    \llap{\raisebox{-\dimexpr1in+\voffset+\topmargin+\headheight+\headsep
      +\textheight+\footskip\relax}{\lycee@licence}}}}}\fi}

\newcommand{\lycee@site}{faire-des-maths.fr}
\newcommand{\lycee@version}{{\scriptsize\color{gray}\lycee@site
  \ifdefined\lyceeversion\ --- Version~\lyceeversion\fi}}
\AtBeginDocument{\iflycee@licence\AddToHookNext{shipout/foreground}{%
  \put(0,0){\rlap{\kern\dimexpr1in+\hoffset+\oddsidemargin\relax
    \raisebox{-\dimexpr1in+\voffset+\topmargin+\headheight+\headsep
      +\textheight+\footskip\relax}{\lycee@version}}}}\fi}

\newcommand{\formulesserrees}{%
  \setlength{\abovedisplayskip}{3pt}\setlength{\belowdisplayskip}{3pt}%
  \setlength{\abovedisplayshortskip}{2pt}\setlength{\belowdisplayshortskip}{3pt}}

\setlength{\parindent}{0pt}

\RequirePackage[implicit=false,hidelinks,pdfencoding=auto,psdextra,bookmarksopen,
  bookmarksopenlevel=1,pdfcreator={},pdfproducer={}]{hyperref}
\RequirePackage{bookmark}
\hypersetup{pdfauthor={Luc Destombe},pdfkeywords={CC BY-NC-SA 4.0}}
\newcounter{lycee@signet}
\newcommand{\signet}[3][]{\stepcounter{lycee@signet}%
  \edef\lycee@dest{\if\relax\detokenize{#1}\relax
    signet.\arabic{lycee@signet}\else#1\fi}%
  \hypertarget{\lycee@dest}{}\bookmark[level=#2,dest=\lycee@dest]{#3}}
\pdfstringdefDisableCommands{%
  \def\R{R}\def\vect#1{#1}\def\Cf{Cf}\def\Cg{Cg}\def\fois#1{×#1}%
  \def\barre#1{#1}\def\intervalle#1#2{[#1;#2]}\def\intff#1#2{[#1;#2]}%
  \def\textsuperscript#1{#1}\def\dfrac#1#2{#1/#2}\def\frac#1#2{#1/#2}%
  \def\vec#1{#1⃗}}%
\RequirePackage[official]{eurosym}

\tcbset{exercice corrige/.style={
  enhanced,breakable,before upper=\raggedright,
  colback=white,colframe=coulExo,
  boxrule=0.8pt,arc=2pt,
  left=6pt,right=6pt,top=8pt,bottom=5pt,
  before skip=9pt,after skip=4pt,
  coltitle=white,fonttitle=\bfseries\small,
  attach boxed title to top left={xshift=8pt,yshift=-\tcboxedtitleheight/2},
  boxed title style={colback=coulExo,arc=2pt,boxrule=0pt}}}
\newcounter{exo}
\newtcolorbox{exercice}[1][]{exercice corrige,
  phantom={\signet[exercice-\arabic{exo}]{2}{Exercice \arabic{exo}}},
  title={Exercice \theexo\ifx\relax#1\relax\else{} --- #1\fi}}
\newenvironment{exo}[1][\relax]{\refstepcounter{exo}\begin{exercice}[#1]}{\end{exercice}}

\geometry{top=1.8cm,bottom=1cm,left=1cm,right=1cm,headheight=14pt,footskip=0.6cm}
\renewcommand{\lycee@pied}{\thepage}

\renewcommand{\theexo}{\ifnum\value{exo}<10 0\fi\arabic{exo}}
\tcbset{exercice corrige/.style={
  enhanced,breakable,before upper=\raggedright,
  colback=white,colframe=coulExo,
  boxrule=0.9pt,arc=2pt,
  left=8pt,right=8pt,top=8pt,bottom=6pt,
  before skip=8pt,after skip=6pt,
  coltitle=white,fonttitle=\bfseries,
  attach boxed title to top left={xshift=8pt,yshift=-\tcboxedtitleheight/2},
  boxed title style={colback=coulExo,arc=2pt,boxrule=0pt}}}

\newcommand{\entetefiche}[2]{%
  \begin{center}
    \begin{tcolorbox}[enhanced,width=0.92\linewidth,colback=white,colframe=coulPart,
        boxrule=1.4pt,arc=3pt,halign=center,top=6pt,bottom=6pt]
      {\large\bfseries\color{coulPart} CORRIGÉ --- FICHE D'EXERCICES}\\[3pt]
      {\normalsize\bfseries #1}\\[2pt]
      {\normalsize #2}
    \end{tcolorbox}
  \end{center}}

\newcounter{partie}
\newcommand{\partie}[1]{%
  \stepcounter{partie}%
  \besoinplace{8\baselineskip}\vspace{4pt}%
  \begin{tcolorbox}[enhanced,colback=coulPart!8,colframe=coulPart,boxrule=0pt,
      leftrule=3pt,arc=0pt,left=6pt,right=4pt,top=3pt,bottom=3pt,
      before skip=6pt,after skip=2pt]
    \bfseries\color{coulPart}\Roman{partie} --- #1\signet{1}{\Roman{partie} --- #1}
  \end{tcolorbox}}
\newcommand{\partiesansnum}[1]{%
  \besoinplace{8\baselineskip}\vspace{4pt}%
  \begin{tcolorbox}[enhanced,colback=coulPart!8,colframe=coulPart,boxrule=0pt,
      leftrule=3pt,arc=0pt,left=6pt,right=4pt,top=3pt,bottom=3pt,
      before skip=6pt,after skip=2pt]
    \bfseries\color{coulPart}#1\signet{1}{#1}
  \end{tcolorbox}}

\setlist[questions,1]{label=\arabic*\textdegree),leftmargin=*,itemsep=3pt,topsep=3pt}
\setlist[questions,2]{label=\alph*),leftmargin=*,itemsep=2pt,topsep=2pt}
\setlist[questions,3]{label=\roman*),leftmargin=*,itemsep=1pt,topsep=1pt}
\newlist{expr}{enumerate}{1}
\setlist[expr]{label={\Alph*\ $=$},leftmargin=*,itemsep=2pt,topsep=3pt}

\newcommand{\res}[1]{{\color{coulRep}#1}}
\newcommand{\rem}[1]{\par\smallskip{\small\itshape\color{black!60}$\rightarrow$ #1}}
\newenvironment{calculs}{\par\smallskip\noindent$\displaystyle\begin{aligned}[t]}%
  {\end{aligned}$\par\smallskip}
\newcommand{\justif}[1]{\quad\text{\small\itshape\color{black!60}#1}}
\newcommand{\cas}[1]{\par\smallskip\noindent\textbf{\color{coulExo}#1}\vadjust{\nobreak}}
\newcommand{\explique}[2]{\par\smallskip\noindent
  \begin{minipage}[t]{0.6\linewidth}\raggedright #1\end{minipage}\hfill
  \begin{minipage}[t]{0.37\linewidth}\raggedright\leavevmode
    {\small\itshape\color{black!60}#2}\end{minipage}\par}

\setlength{\parskip}{2pt}

\makeatother
\entete{Chapitre 8 --- Repérage, vecteurs et trigonométrie --- corrigé}{Seconde}
\setlength{\parindent}{0pt}
\setlength{\jot}{5pt}

\let\casligne\cas
\renewcommand{\cas}[1]{\casligne{#1}\par\nopagebreak}

\ShorthandsOff
\tikzset{
  vecteur/.style={-{Stealth[length=2.2mm]},line width=0.9pt,coulCourbe},
  vecteurB/.style={-{Stealth[length=2.2mm]},line width=0.9pt,coulCourbeB},
  vecteurC/.style={-{Stealth[length=2.2mm]},line width=0.9pt,coulCourbeC},
  cherche/.style={-{Stealth[length=2.4mm]},line width=1.1pt,black},
  base/.style={-{Stealth[length=2mm]},line width=0.8pt,black},
  aide/.style={dashed,black!50,line width=0.5pt},
  chemin/.style={dashed,->,black!60,line width=0.7pt},
  fig/.style={line width=0.7pt,black!75},
  nom/.style={midway,font=\footnotesize,inner sep=1.5pt},
}
\ShorthandsOn
\newcommand{\coord}[2]{\begin{pmatrix}#1\\#2\end{pmatrix}}

\newcommand{\cotefig}[3][0.32]{\par\vspace{-0.55\baselineskip}\noindent
  \begin{minipage}[t]{\dimexpr0.97\linewidth-#1\linewidth\relax}\raggedright #2\end{minipage}\hfill
  \begin{minipage}[t]{#1\linewidth}\vspace{0pt}\centering #3\end{minipage}\par}

\begin{document}

\entetefiche{Chapitre 8 --- Repérage, vecteurs et trigonométrie}{Seconde}

{\small Le repère $(O\,;\vec{\imath},\vec{\jmath})$ est orthonormé. Coordonnées d'un
vecteur : le déplacement, horizontal puis vertical, de l'origine à la pointe ; par le
calcul, \og extrémité moins origine\fg.\par}

\partie{Repère du plan, base orthonormée}

\begin{exo}[Lire et placer]
\cas{a)}
\explique{$\res{A(2\,;3)}$, $\res{B(-3\,;2)}$, $\res{C(-4\,;-2)}$ et $\res{D(1\,;-3)}$.}{d'abord
l'abscisse (axe horizontal), puis l'ordonnée}
\explique{$\res{E(0\,;2)}$.}{$E$ est sur l'axe des ordonnées : son abscisse est $0$}
\cas{b)}
\cotefig{%
\explique{$F$ : $3$ carreaux vers la droite depuis $O$, puis $1$ vers le bas.}{abscisse $3$,
ordonnée $-1$}
\explique{$G$ : $2$ carreaux vers la gauche, sans monter : $G$ est sur l'axe des
abscisses.}{ordonnée $0$}
\explique{$H$ : on ne se déplace pas horizontalement, puis $4$ carreaux vers le bas : $H$
est sur l'axe des ordonnées.}{abscisse $0$}}
{\begin{tikzpicture}[x=0.36cm,y=0.36cm]
  \repere{-3}{-5}{4}{1}
  \gradx{-2,3} \grady{-4,-1}
  \pt{3,-1}{right}{F} \pt{-2,0}{above}{G} \pt{0,-4}{right}{H}
\end{tikzpicture}}
\cas{c)}
\explique{$\vect{OA}=\res{2\vec{\imath}+3\vec{\jmath}}$}{de $O$ à $A$ : $2$ fois
$\vec{\imath}$, puis $3$ fois $\vec{\jmath}$ ; les coefficients sont les coordonnées de $A$}
\explique{$\vect{OC}=\res{-4\vec{\imath}-2\vec{\jmath}}$}{$4$ pas vers la gauche, $2$ vers
le bas : coefficients négatifs}
\cas{d)}
\explique{$\vect{CA}=\res{6\vec{\imath}+5\vec{\jmath}}$}{de $C(-4\,;-2)$ à $A(2\,;3)$ :
$6$ carreaux vers la droite, puis $5$ vers le haut}
\rem{Pour un vecteur qui ne part pas de $O$, on compte le déplacement entre ses deux
points, pas les coordonnées de l'extrémité.}
\end{exo}

\begin{exo}[Abscisse et ordonnée]
\explique{Le point $K(-2\,;5)$ a pour abscisse \res{$-2$} et pour ordonnée \res{$5$} ; il
n'est sur \res{aucun axe}.}{aucune de ses coordonnées n'est nulle}
\explique{Le point $L(0\,;3)$ a pour abscisse \res{$0$} et pour ordonnée \res{$3$} ; il est
sur \res{l'axe des ordonnées}.}{abscisse nulle : on ne se déplace pas horizontalement}
\explique{Le point $M(4\,;0)$ a pour abscisse \res{$4$} et pour ordonnée \res{$0$} ; il est
sur \res{l'axe des abscisses}.}{ordonnée nulle : on ne monte pas}
\rem{Piège : le point d'abscisse nulle est sur l'axe des \emph{ordonnées}, et
réciproquement.}
\end{exo}

\begin{exo}[Points définis par $\vec{\imath}$ et $\vec{\jmath}$]
\cotefig{%
\explique{$\vect{OR}=3\vec{\imath}-2\vec{\jmath}$, donc $\res{R(3\,;-2)}$.}{les
coefficients de $\vec{\imath}$ et de $\vec{\jmath}$ sont les coordonnées (définition 2)}
\explique{$\vect{OS}=-\vec{\imath}+4\vec{\jmath}$, donc $\res{S(-1\,;4)}$.}{$-\vec{\imath}$ :
coefficient $-1$}
\explique{$\vect{OT}=-2\vec{\jmath}$, soit $0\vec{\imath}-2\vec{\jmath}$, donc
$\res{T(0\,;-2)}$.}{pas de $\vec{\imath}$ : abscisse $0$}
\explique{On place chaque point en partant de $O$ : d'abord les pas de $\vec{\imath}$, puis
ceux de $\vec{\jmath}$.}{figure ci-contre}}
{\begin{tikzpicture}[x=0.36cm,y=0.36cm]
  \repere{-2}{-3}{4}{5}
  \gradx{-1,3} \grady{-2,4}
  \pt{3,-2}{right}{R} \pt{-1,4}{left}{S} \pt{0,-2}{left}{T}
\end{tikzpicture}}
\end{exo}

\begin{exo}[Le repère d'un parallélogramme]
\cas{a)}
\cotefig[0.36]{%
\explique{On trace un parallélogramme quelconque $ABCD$ et son centre $K$, intersection des
diagonales.}{$K$ est le milieu de $[AC]$ et de $[BD]$}}
{\begin{tikzpicture}[scale=0.55]
  \draw[fig] (0,0) -- (4,0) -- (5.5,2.5) -- (1.5,2.5) -- cycle;
  \draw[black!40] (0,0) -- (5.5,2.5) (4,0) -- (1.5,2.5);
  \draw[vecteur] (0,0) -- node[nom,below]{$\vect{AB}$} (4,0);
  \draw[vecteurB] (0,0) -- node[nom,left]{$\vect{AD}$} (1.5,2.5);
  \pt{0,0}{below left}{A} \pt{4,0}{below right}{B} \pt{5.5,2.5}{above right}{C}
  \pt{1.5,2.5}{above left}{D} \pt{2.75,1.25}{below}{K}
\end{tikzpicture}}
\cas{b)}
\explique{$\vect{AA}=0\vect{AB}+0\vect{AD}$, donc $\res{A(0\,;0)}$ : c'est l'origine.}{on
écrit chaque vecteur $\vect{AM}$ avec $\vect{AB}$ et $\vect{AD}$}
\explique{$\vect{AB}=1\vect{AB}+0\vect{AD}$, donc $\res{B(1\,;0)}$ ; de même
$\res{D(0\,;1)}$.}{$\vect{AB}$ et $\vect{AD}$ jouent le rôle de $\vec{\imath}$ et
$\vec{\jmath}$}
\cas{c)}
\begin{calculs}
\vect{AC} &= \vect{AB}+\vect{BC} &&\justif{Chasles : on passe par $B$}\\
\vect{AC} &= \res{\vect{AB}+\vect{AD}} &&\justif{$ABCD$ parallélogramme : $\vect{BC}=\vect{AD}$}
\end{calculs}
\explique{Donc $\res{C(1\,;1)}$.}{coefficients $1$ et $1$}
\cas{d)}
\begin{calculs}
\vect{AK} &= \tfrac{1}{2}\vect{AC} &&\justif{$K$ milieu de $[AC]$}\\
\vect{AK} &= \tfrac{1}{2}\vect{AB}+\tfrac{1}{2}\vect{AD} &&\justif{on remplace $\vect{AC}$ (question c) et on développe}
\end{calculs}
\explique{Donc $\res{K\left(\frac{1}{2}\,;\frac{1}{2}\right)}$.}{le repère n'est pas
orthonormé, mais les coordonnées se lisent de la même façon}
\end{exo}

\partie{Coordonnées d'un vecteur}

\begin{exo}[Lire des coordonnées]
\cas{a)}
\explique{$\res{\vec{a}\coord{4}{-1}}$}{$4$ carreaux vers la droite, $1$ vers le bas}
\explique{$\res{\vec{b}\coord{-3}{-3}}$}{$3$ vers la gauche, $3$ vers le bas}
\explique{$\res{\vec{c}\coord{0}{4}}$}{pas de déplacement horizontal : abscisse $0$}
\explique{$\res{\vec{d}\coord{-4}{1}}$}{$4$ vers la gauche, $1$ vers le haut}
\explique{$\res{\vec{e}\coord{3}{0}}$}{pas de déplacement vertical : ordonnée $0$}
\cas{b)}
\explique{$\res{\vec{a}}$ et $\res{\vec{d}}$ sont opposés : leurs coordonnées sont
opposées, $\coord{4}{-1}$ et $\coord{-4}{1}$.}{l'opposé de $\coord{x}{y}$ est
$\coord{-x}{-y}$}
\end{exo}

\begin{exo}[Tracer des représentants]
\cas{a)}
\cotefig[0.4]{%
\explique{Depuis $A(-1\,;1)$, on fait chaque déplacement et on place la pointe.}{la flèche
part de $A$}
\explique{$\vec{u}$ : pointe en $\res{(2\,;3)}$ ; $\vec{v}$ : en $\res{(-3\,;4)}$.}{$3$ à
droite, $2$ en haut ; $2$ à gauche, $3$ en haut}
\explique{$\vec{w}$ : en $\res{(-1\,;-2)}$ ; $\vec{t}$ : en $\res{(-5\,;1)}$.}{$3$ vers le
bas ; $4$ vers la gauche}}
{\begin{tikzpicture}[x=0.36cm,y=0.36cm]
  \repere{-6}{-3}{3}{5}
  \gradx{-5,-3,-1,2} \grady{-2,1,3}
  \draw[cherche] (-1,1) -- node[nom,below right]{$\vec{u}$} (2,3);
  \draw[cherche] (-1,1) -- node[nom,above right]{$\vec{v}$} (-3,4);
  \draw[cherche] (-1,1) -- node[nom,right]{$\vec{w}$} (-1,-2);
  \draw[cherche] (-1,1) -- node[nom,below]{$\vec{t}$} (-5,1);
  \pt{-1,1}{above left}{A}
\end{tikzpicture}}
\cas{b)}
\cotefig[0.4]{%
\explique{La pointe doit être en $B(4\,;-2)$ : depuis $B$, on fait le déplacement à
l'envers, $3$ vers la gauche et $1$ vers le haut. L'origine est $\res{(1\,;-1)}$.}{$B$ est
l'extrémité}}
{\begin{tikzpicture}[x=0.36cm,y=0.36cm]
  \repere{-1}{-3}{5}{1}
  \gradx{1,4} \grady{-2}
  \draw[cherche] (1,-1) -- node[nom,above right]{$\vec{s}$} (4,-2);
  \pt{4,-2}{below}{B}
\end{tikzpicture}}
\end{exo}

\begin{exo}[Calculer des coordonnées]
\explique{Pour chaque vecteur, extrémité moins origine (propriété 1).}{abscisse avec les
abscisses, ordonnée avec les ordonnées}
\begin{calculs}
\vect{AB} &\coord{6-2}{1-5}=\res{\coord{4}{-4}} &&\justif{$A(2\,;5)$, $B(6\,;1)$}\\
\vect{CD} &\coord{1-(-3)}{-2-4}=\res{\coord{4}{-6}} &&\justif{parenthèses autour de $-3$}\\
\vect{DE} &\coord{-4-1}{-1-(-2)}=\res{\coord{-5}{1}} &&\justif{$-1+2=1$}\\
\vect{BC} &\coord{-3-6}{4-1}=\res{\coord{-9}{3}} &&\justif{origine $B$, extrémité $C$}\\
\vect{EA} &\coord{2-(-4)}{5-(-1)}=\res{\coord{6}{6}} &&\justif{origine $E$, extrémité $A$}
\end{calculs}
\explique{$\vect{AB}=\res{4\vec{\imath}-4\vec{\jmath}}$.}{les coordonnées sont les
coefficients de $\vec{\imath}$ et $\vec{\jmath}$}
\end{exo}

\begin{exo}[Lire, puis calculer]
\cas{a)}
\explique{$\res{A(-3\,;1)}$, $\res{B(1\,;3)}$, $\res{C(3\,;-1)}$ et $\res{D(-1\,;-3)}$.}{abscisse,
puis ordonnée}
\cas{b)}
\begin{calculs}
\vect{AB} &\coord{1-(-3)}{3-1}=\res{\coord{4}{2}} &&\justif{extrémité moins origine}\\
\vect{DC} &\coord{3-(-1)}{-1-(-3)}=\res{\coord{4}{2}} &&\justif{origine $D$, extrémité $C$}\\
\vect{BC} &\coord{3-1}{-1-3}=\res{\coord{2}{-4}} &&
\end{calculs}
\cas{c)}
\explique{De $A$ à $B$ : $4$ carreaux vers la droite, $2$ vers le haut ; de $D$ à $C$ :
pareil ; de $B$ à $C$ : $2$ vers la droite, $4$ vers le bas.}{le comptage confirme le
calcul}
\cas{d)}
\explique{$\vect{AB}$ et $\vect{DC}$ ont les mêmes coordonnées : $\res{\vect{AB}=\vect{DC}}$.}{donc
$ABCD$ est un parallélogramme (partie V)}
\end{exo}

\begin{exo}[Un point inconnu]
\cas{a)}
\explique{On note $E(x\,;y)$ : $\vect{AE}\coord{x-3}{y-2}$ doit valoir
$\coord{-1}{4}$.}{extrémité moins origine}
\begin{calculs}
x-3 &= -1 &&\justif{abscisses}\\
x &= 2 &&\justif{on ajoute $3$}
\end{calculs}
\begin{calculs}
y-2 &= 4 &&\justif{ordonnées}\\
y &= 6 &&\justif{on ajoute $2$}
\end{calculs}
\explique{Donc $\res{E(2\,;6)}$.}{contrôle : de $A(3\,;2)$ à $E(2\,;6)$, $1$ à gauche,
$4$ en haut}
\cas{b)}
\explique{On note $F(x\,;y)$ : $\vect{FB}\coord{-5-x}{9-y}$ doit valoir
$\coord{2}{8}$.}{cette fois, $F$ est l'origine}
\begin{calculs}
-5-x &= 2 &&\justif{abscisses}\\
-x &= 7 &&\justif{on ajoute $5$}\\
x &= -7 &&
\end{calculs}
\begin{calculs}
9-y &= 8 &&\justif{ordonnées}\\
y &= 1 &&
\end{calculs}
\explique{Donc $\res{F(-7\,;1)}$.}{contrôle : $-5-(-7)=2$ et $9-1=8$}
\rem{Erreur fréquente : écrire $x_{F}-x_{B}$ ; l'origine de $\vect{FB}$ est $F$, c'est son
abscisse qu'on retranche.}
\end{exo}

\partie{Coordonnées d'une somme, d'un produit par un nombre réel}

\begin{exo}[Calculs de coordonnées]
\cas{a)}
\begin{calculs}
\vec{u}+\vec{v} &\coord{2+(-4)}{-3+5}=\res{\coord{-2}{2}} &&\justif{abscisses ensemble, ordonnées ensemble}
\end{calculs}
\cas{b)}
\begin{calculs}
3\vec{w} &\coord{3\times 3}{3\times 2}=\res{\coord{9}{6}} &&\justif{on multiplie les deux coordonnées par $3$}
\end{calculs}
\cas{c)}
\begin{calculs}
\vec{u}-2\vec{v} &\coord{2-2\times(-4)}{-3-2\times 5}=\res{\coord{10}{-13}} &&\justif{$2+8=10$ ; $-3-10=-13$}
\end{calculs}
\cas{d)}
\begin{calculs}
-0{,}5\vec{v} &\coord{-0{,}5\times(-4)}{-0{,}5\times 5}=\res{\coord{2}{-2{,}5}} &&\justif{on multiplie par $-0{,}5$}
\end{calculs}
\cas{e)}
\begin{calculs}
x &= 2\times 2+(-4)-3 &&\justif{abscisse : $2x_{u}+x_{v}-x_{w}$}\\
x &= -3 &&\\
y &= 2\times(-3)+5-2 &&\justif{ordonnée : mêmes opérations}\\
y &= -3 &&
\end{calculs}
\explique{Donc $2\vec{u}+\vec{v}-\vec{w}$ a pour coordonnées $\res{\coord{-3}{-3}}$.}{on
conclut}
\end{exo}

\begin{exo}[Avec des points]
\cas{a)}
\begin{calculs}
\vect{AB} &\coord{4-1}{-1-2}=\res{\coord{3}{-3}} &&\justif{extrémité moins origine}\\
\vect{AC} &\coord{-2-1}{3-2}=\res{\coord{-3}{1}} &&
\end{calculs}
\cas{b)}
\begin{calculs}
\vect{AB}+\vect{AC} &\coord{3+(-3)}{-3+1}=\res{\coord{0}{-2}} &&\justif{somme}\\
2\vect{AB}-\vect{AC} &\coord{6-(-3)}{-6-1}=\res{\coord{9}{-7}} &&\justif{$2\vect{AB}\coord{6}{-6}$}
\end{calculs}
\cas{c)}
\begin{calculs}
\vect{BC} &\coord{-2-4}{3-(-1)}=\res{\coord{-6}{4}} &&\justif{extrémité moins origine}\\
\vect{AC}-\vect{AB} &\coord{-3-3}{1-(-3)}=\coord{-6}{4} &&\justif{différence}
\end{calculs}
\explique{Mêmes coordonnées : $\res{\vect{BC}=\vect{AC}-\vect{AB}}$.}{c'est la relation de
Chasles : $\vect{BC}=\vect{BA}+\vect{AC}$}
\end{exo}

\begin{exo}[Un point défini par une somme]
\cas{a)}
\begin{calculs}
\vec{u}+\vec{v} &\coord{2+(-1)}{-3+5}=\res{\coord{1}{2}} &&\justif{somme}\\
2\vec{u}-\vec{v} &\coord{4-(-1)}{-6-5}=\res{\coord{5}{-11}} &&\justif{$2\vec{u}\coord{4}{-6}$}
\end{calculs}
\cas{b)}
\explique{$\vect{AE}=\vec{u}+\vec{v}$ : depuis $A(-1\,;3)$, on ajoute $1$ à l'abscisse et
$2$ à l'ordonnée, donc $\res{E(0\,;5)}$.}{$x_{E}-(-1)=1$ et $y_{E}-3=2$}
\explique{$\vect{AF}=2\vec{u}-\vec{v}$ : $x_{F}=-1+5=4$ et $y_{F}=3-11=-8$, donc
$\res{F(4\,;-8)}$.}{même raisonnement}
\cas{c)}
\begin{calculs}
\vect{EF} &\coord{4-0}{-8-5}=\res{\coord{4}{-13}} &&\justif{extrémité moins origine (partie II)}
\end{calculs}
\end{exo}

\begin{exo}[Démontrer une égalité]
\cas{a)}
\begin{calculs}
2\vec{u} &\coord{2\times 2}{2\times 1}=\coord{4}{2} &&\justif{produit}\\
\vec{v}+\vec{w} &\coord{1+3}{-2+4}=\coord{4}{2} &&\justif{somme}
\end{calculs}
\explique{Mêmes coordonnées, donc $\res{2\vec{u}=\vec{v}+\vec{w}}$.}{deux vecteurs égaux ont
les mêmes coordonnées}
\cas{b)}
\cotefig[0.36]{%
\explique{Depuis un même point, $2\vec{u}$ arrive au même endroit que $\vec{v}$ suivi de
$\vec{w}$ : \res{les deux pointes coïncident}.}{c'est l'égalité du a), vue sur le
quadrillage}}
{\begin{tikzpicture}[x=0.38cm,y=0.38cm]
  \quadrillage{-0.4}{-2.4}{4.4}{2.4}
  \draw[cherche] (0,0) -- node[nom,above left]{$2\vec{u}$} (4,2);
  \draw[vecteur] (0,0) -- node[nom,below left]{$\vec{v}$} (1,-2);
  \draw[vecteurB] (1,-2) -- node[nom,right]{$\vec{w}$} (4,2);
\end{tikzpicture}}
\end{exo}

\partie{Combinaison de deux vecteurs}

\begin{exo}[Lire une combinaison]
\explique{Pour chaque vecteur, on cherche un chemin de son origine à sa pointe fait de
$\vec{a}\coord{3}{1}$ et de $\vec{b}\coord{1}{3}$ (ou de leurs opposés).}{méthode 5}
\cotefig[0.41]{%
\explique{$\vec{w}_{1}$ : $\vec{a}$ puis $\vec{b}$, donc
$\vec{w}_{1}=\res{\vec{a}+\vec{b}}$.}{contrôle : $\coord{3}{1}+\coord{1}{3}=\coord{4}{4}$}
\explique{$\vec{w}_{2}$ : deux fois $\vec{a}$, puis $-\vec{b}$, donc
$\vec{w}_{2}=\res{2\vec{a}-\vec{b}}$.}{$\coord{6}{2}-\coord{1}{3}=\coord{5}{-1}$}
\explique{$\vec{w}_{3}$ : $-\vec{a}$, puis deux fois $\vec{b}$, donc
$\vec{w}_{3}=\res{-\vec{a}+2\vec{b}}$.}{$\coord{-3}{-1}+\coord{2}{6}=\coord{-1}{5}$}}
{\begin{tikzpicture}[x=0.32cm,y=0.32cm]
  \quadrillage{-0.4}{-1.4}{18.4}{6.4}
  \draw[cherche] (4,0) -- node[nom,above left]{$\vec{w}_{1}$} (8,4);
  \draw[chemin] (4,0) -- node[nom,below right]{$\vec{a}$} (7,1);
  \draw[chemin] (7,1) -- node[nom,right]{$\vec{b}$} (8,4);
  \draw[cherche] (9,3) -- node[nom,below]{$\vec{w}_{2}$} (14,2);
  \draw[chemin] (9,3) -- node[nom,above left]{$\vec{a}$} (12,4);
  \draw[chemin] (12,4) -- node[nom,above left]{$\vec{a}$} (15,5);
  \draw[chemin] (15,5) -- node[nom,right]{$-\vec{b}$} (14,2);
  \draw[cherche] (17,0) -- node[nom,right]{$\vec{w}_{3}$} (16,5);
  \draw[chemin] (17,0) -- node[nom,below]{$-\vec{a}$} (14,-1);
  \draw[chemin] (14,-1) -- node[nom,left]{$\vec{b}$} (15,2);
  \draw[chemin] (15,2) -- node[nom,left]{$\vec{b}$} (16,5);
\end{tikzpicture}}
\end{exo}

\begin{exo}[Construire une combinaison]
\explique{$\vec{c}$ fait $4$ carreaux vers la droite et $1$ vers le haut ; $\vec{d}$ fait
$1$ vers la gauche et $2$ vers le haut.}{on lit d'abord les deux vecteurs}
\cotefig[0.45]{%
\explique{$\vec{r}=\vec{c}+2\vec{d}$ : depuis $M$, on reporte $\vec{c}$, puis deux fois
$\vec{d}$ ; la pointe de $\vec{r}$ est au bout du chemin, $\res{2}$ carreaux à droite et
$\res{5}$ vers le haut de $M$.}{$\vec{r}\coord{2}{5}$}
\explique{$\vec{s}=2\vec{c}-\vec{d}$ : deux fois $\vec{c}$, puis $-\vec{d}$ ($1$ à droite,
$2$ vers le bas) ; la pointe est $\res{9}$ carreaux à droite de $M$, au même niveau.}{$\vec{s}\coord{9}{0}$}}
{\begin{tikzpicture}[x=0.4cm,y=0.4cm]
  \quadrillage{-0.4}{-0.4}{14.4}{6.4}
  \draw[chemin] (3,0) -- node[nom,below right]{$\vec{c}$} (7,1);
  \draw[chemin] (7,1) -- node[nom,right]{$\vec{d}$} (6,3);
  \draw[chemin] (6,3) -- node[nom,right]{$\vec{d}$} (5,5);
  \draw[cherche] (3,0) -- node[nom,left]{$\vec{r}$} (5,5);
  \draw[chemin] (7,1) -- node[nom,above left]{$\vec{c}$} (11,2);
  \draw[chemin] (11,2) -- node[nom,right]{$-\vec{d}$} (12,0);
  \draw[cherche] (3,0) -- node[nom,below]{$\vec{s}$} (12,0);
  \pt{3,0}{below left}{M}
\end{tikzpicture}}
\end{exo}

\begin{exo}[Dans un parallélogramme, dans un triangle]
\cas{a)}
\begin{calculs}
\vect{BD} &= \vect{BA}+\vect{AD} &&\justif{Chasles : on passe par $A$}\\
\vect{BD} &= \res{-\vect{AB}+\vect{AD}} &&\justif{$\vect{BA}=-\vect{AB}$}
\end{calculs}
\begin{calculs}
\vect{BK} &= \tfrac{1}{2}\vect{BD} &&\justif{$K$ milieu de la diagonale $[BD]$}\\
\vect{BK} &= \tfrac{1}{2}\bigl(-\vect{AB}+\vect{AD}\bigr) &&\justif{on remplace $\vect{BD}$}\\
\vect{BK} &= \res{-\tfrac{1}{2}\vect{AB}+\tfrac{1}{2}\vect{AD}} &&\justif{on développe}
\end{calculs}
\cas{b)}
\begin{calculs}
\vect{EM} &= \vect{EF}+\vect{FM} &&\justif{Chasles : on passe par $F$}\\
\vect{EM} &= \vect{EF}+\tfrac{2}{3}\vect{FG} &&\justif{énoncé}\\
\vect{EM} &= \vect{EF}+\tfrac{2}{3}\bigl(\vect{FE}+\vect{EG}\bigr) &&\justif{Chasles : on passe par $E$}\\
\vect{EM} &= \vect{EF}-\tfrac{2}{3}\vect{EF}+\tfrac{2}{3}\vect{EG} &&\justif{$\vect{FE}=-\vect{EF}$}\\
\vect{EM} &= \res{\tfrac{1}{3}\vect{EF}+\tfrac{2}{3}\vect{EG}} &&\justif{$1-\frac23=\frac13$}
\end{calculs}
\cas{c)}
\explique{Dans le repère $(E\,;\vect{EF},\vect{EG})$, les coordonnées de $M$ sont les
coefficients de $\vect{EM}$ : $\res{M\left(\frac{1}{3}\,;\frac{2}{3}\right)}$.}{comme
$\vect{OM}=x\vec{\imath}+y\vec{\jmath}$ (définition 2)}
\end{exo}

\begin{exo}[Combinaisons et coordonnées]
\cas{a)}
\begin{calculs}
2\vec{a}+\vec{b} &\coord{2\times 2+(-1)}{2\times 1+3}=\res{\coord{3}{5}} &&\justif{$2\vec{a}\coord{4}{2}$}\\
\vec{a}-\vec{b} &\coord{2-(-1)}{1-3}=\res{\coord{3}{-2}} &&\justif{différence}
\end{calculs}
\cas{b)}
\begin{calculs}
\vec{a}+2\vec{b} &\coord{2+2\times(-1)}{1+2\times 3}=\coord{0}{7} &&\justif{$2\vec{b}\coord{-2}{6}$}
\end{calculs}
\explique{Ce sont les coordonnées de $\vec{w}$, donc $\res{\vec{w}=\vec{a}+2\vec{b}}$.}{mêmes
coordonnées}
\cas{c)}
\cotefig[0.3]{%
\explique{De l'origine de $\vec{w}$ : $\vec{a}$ ($2$ à droite, $1$ en haut), puis deux fois
$\vec{b}$ ($1$ à gauche, $3$ en haut) ; on arrive \res{à la pointe de $\vec{w}$}.}{le
chemin redonne $\coord{0}{7}$}}
{\begin{tikzpicture}[x=0.32cm,y=0.32cm]
  \quadrillage{-0.4}{-0.4}{3.4}{7.4}
  \draw[cherche] (1,0) -- node[nom,left]{$\vec{w}$} (1,7);
  \draw[vecteur] (1,0) -- node[nom,right]{$\vec{a}$} (3,1);
  \draw[vecteurB] (3,1) -- node[nom,right]{$\vec{b}$} (2,4);
  \draw[vecteurB] (2,4) -- node[nom,right]{$\vec{b}$} (1,7);
\end{tikzpicture}}
\cas{d)}
\begin{calculs}
\vec{a}-2\vec{b} &\coord{2-2\times(-1)}{1-2\times 3}=\coord{4}{-5} &&\justif{$2+2=4$ ; $1-6=-5$}
\end{calculs}
\explique{$x_{M}=1+4=5$ et $y_{M}=1-5=-4$, donc $\res{M(5\,;-4)}$.}{$\vect{AM}\coord{x_{M}-1}{y_{M}-1}$}
\end{exo}

\partie{Égalité de vecteurs et parallélogrammes}

\begin{exo}[Parallélogramme ou non]
\explique{$ABCD$ est un parallélogramme quand $\vect{AB}=\vect{DC}$ : on compare les
coordonnées.}{propriété 4 ; ordre des lettres}
\cas{a)}
\begin{calculs}
\vect{AB} &\coord{1-(-2)}{3-1}=\coord{3}{2} &&\justif{extrémité moins origine}\\
\vect{DC} &\coord{4-1}{1-(-1)}=\coord{3}{2} &&\justif{origine $D$, extrémité $C$}
\end{calculs}
\explique{$\vect{AB}=\vect{DC}$ : $ABCD$ \res{est un parallélogramme}.}{mêmes coordonnées}
\cas{b)}
\begin{calculs}
\vect{KL} &\coord{2-(-1)}{-1-(-2)}=\coord{3}{1} &&\\
\vect{NM} &\coord{4-1}{3-2}=\coord{3}{1} &&
\end{calculs}
\explique{$\vect{KL}=\vect{NM}$ : $KLMN$ \res{est un parallélogramme}.}{mêmes coordonnées}
\cas{c)}
\begin{calculs}
\vect{RS} &\coord{3-0}{2-1}=\coord{3}{1} &&\\
\vect{UT} &\coord{4-0}{5-4}=\coord{4}{1} &&
\end{calculs}
\explique{$\vect{RS}\neq\vect{UT}$ : $RSTU$ \res{n'est pas un parallélogramme}.}{les
abscisses diffèrent}
\end{exo}

\begin{exo}[Le quatrième sommet]
\cas{a)}
\explique{$ABCD$ parallélogramme : $\vect{AB}=\vect{DC}$, avec $\vect{AB}\coord{2}{3}$ et,
pour $D(x\,;y)$, $\vect{DC}\coord{3-x}{2-y}$.}{$\vect{AB}\coord{0-(-2)}{4-1}$}
\begin{calculs}
3-x &= 2 &&\justif{abscisses}\\
x &= 1 &&
\end{calculs}
\begin{calculs}
2-y &= 3 &&\justif{ordonnées}\\
y &= -1 &&
\end{calculs}
\explique{Donc $\res{D(1\,;-1)}$.}{contrôle : $\vect{DC}\coord{2}{3}$}
\cas{b)}
\explique{$ABEC$ parallélogramme : $\vect{AB}=\vect{CE}$. Pour $E(x\,;y)$ :
$x-3=2$ et $y-2=3$.}{dans $ABEC$, $[AB]$ est parallèle à $[CE]$}
\explique{Donc $\res{E(5\,;5)}$.}{$x=5$ et $y=5$}
\cas{c)}
\cotefig[0.34]{%
\explique{On place les cinq points et on trace $ABCD$, puis $ABEC$ : ils ont le côté
$[AB]$ en commun.}{figure ci-contre}}
{\begin{tikzpicture}[x=0.32cm,y=0.32cm]
  \repere{-3}{-2}{6}{6}
  \gradx{-2,1,3,5} \grady{-1,1,2,4,5}
  \draw[fig] (-2,1) -- (0,4) -- (3,2) -- (1,-1) -- cycle;
  \draw[fig,coulCourbe] (-2,1) -- (0,4) -- (5,5) -- (3,2) -- cycle;
  \pt{-2,1}{left}{A} \pt{0,4}{above left}{B} \pt{3,2}{right}{C}
  \pt{1,-1}{below}{D} \pt{5,5}{above}{E}
\end{tikzpicture}}
\end{exo}

\begin{exo}[Image par une translation]
\cas{a)}
\explique{$F$ est l'image de $B$ par la translation de vecteur $\vect{AC}$ :
$\res{\vect{BF}=\vect{AC}}$.}{chapitre 3, définition de la translation}
\explique{$\vect{AC}\coord{3}{3}$, donc $x_{F}=2+3=5$ et $y_{F}=-3+3=0$ :
$\res{F(5\,;0)}$.}{$\vect{AC}\coord{4-1}{2-(-1)}$}
\cas{b)}
\begin{calculs}
\vect{AB} &\coord{2-1}{-3-(-1)}=\coord{1}{-2} &&\justif{côté $[AB]$}\\
\vect{CF} &\coord{5-4}{0-2}=\coord{1}{-2} &&\justif{côté opposé $[CF]$, même sens}
\end{calculs}
\explique{$\vect{AB}=\vect{CF}$ : $ABFC$ \res{est un parallélogramme}.}{propriété 4 avec
les lettres $A$, $B$, $F$, $C$}
\cas{c)}
\explique{\res{$C$ est l'image de $A$ par la translation de vecteur $\vect{BF}$}.}{car
$\vect{AC}=\vect{BF}$}
\cas{d)}
\explique{$\vect{AB}+\vect{AC}$ a pour coordonnées $\coord{1+3}{-2+3}=\coord{4}{1}$, donc
$G(1+4\,;-1+1)$, soit $\res{G(5\,;0)}$.}{partie III}
\explique{On remarque que \res{$G=F$} : dans le parallélogramme $ABFC$,
$\vect{AB}+\vect{AC}=\vect{AF}$.}{règle du parallélogramme}
\end{exo}

\partie{Colinéarité, déterminant}

\begin{exo}[Colinéaires ou non]
\explique{Pour chaque couple, on calcule $\det(\vec{u}\,;\vec{v})=xy'-yx'$.}{méthode 7 :
nul, colinéaires ; non nul, pas colinéaires}
\cas{a)}
\begin{calculs}
\det(\vec{u}\,;\vec{v}) &= 4\times 9-(-6)\times(-6) &&\justif{$xy'-yx'$}\\
\det(\vec{u}\,;\vec{v}) &= 36-36 &&\\
\det(\vec{u}\,;\vec{v}) &= 0 &&
\end{calculs}
\explique{\res{Colinéaires}.}{en effet, $\vec{v}=-1{,}5\vec{u}$}
\cas{b)}
\begin{calculs}
\det(\vec{u}\,;\vec{v}) &= 2\times 2-5\times 5 &&\justif{$xy'-yx'$}\\
\det(\vec{u}\,;\vec{v}) &= -21 &&
\end{calculs}
\explique{\res{Pas colinéaires}.}{$-21\neq 0$}
\cas{c)}
\begin{calculs}
\det(\vec{u}\,;\vec{v}) &= \frac{2}{3}\times 3-1\times 2 &&\justif{$xy'-yx'$}\\
\det(\vec{u}\,;\vec{v}) &= 0 &&\justif{$\frac23\times 3=2$}
\end{calculs}
\explique{\res{Colinéaires}.}{$\vec{v}=3\vec{u}$}
\cas{d)}
\begin{calculs}
\det(\vec{u}\,;\vec{v}) &= 3\times 2-(-1)\times(-5) &&\justif{$xy'-yx'$}\\
\det(\vec{u}\,;\vec{v}) &= 1 &&\justif{$6-5$}
\end{calculs}
\explique{\res{Pas colinéaires}.}{$1\neq 0$ : presque, mais pas tout à fait}
\end{exo}

\begin{exo}[Trouver une coordonnée]
\cas{a)}
\explique{Par exemple $\res{\coord{-1}{2}}$ et $\res{\coord{2}{-4}}$.}{$\frac12\vec{u}$ et
$-\vec{u}$ : on multiplie les deux coordonnées par un même nombre}
\cas{b)}
\begin{calculs}
\det(\vec{u}\,;\vec{v}) &= -2\times x-4\times 3 &&\justif{$\vec{u}\coord{-2}{4}$, $\vec{v}\coord{3}{x}$}\\
\det(\vec{u}\,;\vec{v}) &= -2x-12 &&
\end{calculs}
\begin{calculs}
-2x-12 &= 0 &&\justif{colinéaires : déterminant nul}\\
-2x &= 12 &&\\
x &= \res{-6} &&\justif{on divise par $-2$}
\end{calculs}
\explique{Contrôle : $\vec{v}\coord{3}{-6}=-1{,}5\vec{u}$.}{coordonnées proportionnelles}
\end{exo}

\begin{exo}[Parallèles ou non]
\cas{a)}
\explique{$\vect{AB}\coord{3}{1}$ et $\vect{CD}\coord{6}{2}$.}{$\coord{4-1}{3-2}$ et
$\coord{4-(-2)}{1-(-1)}$}
\begin{calculs}
\det(\vect{AB}\,;\vect{CD}) &= 3\times 2-1\times 6 &&\justif{$xy'-yx'$}\\
\det(\vect{AB}\,;\vect{CD}) &= 0 &&
\end{calculs}
\explique{Les droites $(AB)$ et $(CD)$ sont \res{parallèles}.}{vecteurs colinéaires}
\explique{$ABDC$ est un parallélogramme quand $\vect{AB}=\vect{CD}$ ; or
$\vect{CD}=2\vect{AB}$ : $ABDC$ \res{n'est pas un parallélogramme} (c'est un
trapèze).}{partie V : il faut des vecteurs égaux, pas seulement colinéaires}
\cas{b)}
\explique{$\vect{EF}\coord{2}{-4}$ et $\vect{GH}\coord{3}{-5}$.}{$\coord{2-0}{-1-3}$ et
$\coord{2-(-1)}{-5-0}$}
\begin{calculs}
\det(\vect{EF}\,;\vect{GH}) &= 2\times(-5)-(-4)\times 3 &&\justif{$xy'-yx'$}\\
\det(\vect{EF}\,;\vect{GH}) &= 2 &&\justif{$-10+12$}
\end{calculs}
\explique{Les droites $(EF)$ et $(GH)$ \res{ne sont pas parallèles}.}{$2\neq 0$}
\end{exo}

\begin{exo}[Alignés ou non]
\explique{Trois points sont alignés quand deux vecteurs qui partent du même point sont
colinéaires.}{propriété 5}
\cas{a)}
\begin{calculs}
\det(\vect{AB}\,;\vect{AC}) &= -2\times 6-(-4)\times 3 &&\justif{$\vect{AB}\coord{-2}{-4}$, $\vect{AC}\coord{3}{6}$}\\
\det(\vect{AB}\,;\vect{AC}) &= 0 &&\justif{$-12+12$}
\end{calculs}
\explique{$A$, $B$ et $C$ sont \res{alignés}.}{vecteurs colinéaires}
\cas{b)}
\begin{calculs}
\det(\vect{DE}\,;\vect{DF}) &= -2\times 9-(-5)\times 4 &&\justif{$\vect{DE}\coord{-2}{-5}$, $\vect{DF}\coord{4}{9}$}\\
\det(\vect{DE}\,;\vect{DF}) &= 2 &&\justif{$-18+20$}
\end{calculs}
\explique{$D$, $E$ et $F$ \res{ne sont pas alignés}.}{$2\neq 0$}
\cas{c)}
\begin{calculs}
\det(\vect{KL}\,;\vect{KG}) &= -6\times(-4)-4\times 6 &&\justif{$\vect{KL}\coord{-6}{4}$, $\vect{KG}\coord{6}{-4}$}\\
\det(\vect{KL}\,;\vect{KG}) &= 0 &&\justif{$24-24$}
\end{calculs}
\explique{\res{Oui}, $G$ est sur la droite $(KL)$.}{$K$, $L$, $G$ alignés}
\cas{d)}
\explique{$\vect{AB}+\vect{AC}$ a pour coordonnées $\coord{-2+3}{-4+6}=\coord{1}{2}$, donc
$\res{M(2\,;5)}$.}{depuis $A(1\,;3)$ : $1+1=2$ et $3+2=5$}
\begin{calculs}
\det(\vect{AB}\,;\vect{AM}) &= -2\times 2-(-4)\times 1 &&\justif{$\vect{AM}\coord{1}{2}$}\\
\det(\vect{AB}\,;\vect{AM}) &= 0 &&
\end{calculs}
\explique{\res{Oui}, $M$ est sur la droite $(AB)$.}{logique : $\vect{AB}$ et $\vect{AC}$ ont
la même direction, leur somme aussi}
\end{exo}

\partie{Milieu d'un segment}

\begin{exo}[Calculer des milieux]
\explique{Le milieu a pour coordonnées la moyenne des abscisses et la moyenne des
ordonnées.}{propriété 6}
\cas{a)}
\explique{$\left(\frac{2+6}{2}\,;\frac{5+(-1)}{2}\right)$, soit $\res{(4\,;2)}$.}{$\frac82$
et $\frac42$}
\cas{b)}
\explique{$\left(\frac{10+(-4)}{2}\,;\frac{1+5}{2}\right)$, soit $\res{(3\,;3)}$.}{$\frac62$ et
$\frac62$}
\cas{c)}
\explique{$\left(\frac{3+(-3)}{2}\,;\frac{-4+4}{2}\right)$, soit $\res{(0\,;0)}$.}{$A$ et $B$
sont symétriques par rapport à $O$}
\cas{d)}
\explique{$\left(\frac{-1+4}{2}\,;\frac{2+7}{2}\right)$, soit $\res{(1{,}5\,;4{,}5)}$.}{$\frac32$
et $\frac92$}
\end{exo}

\begin{exo}[Retrouver une extrémité]
\cas{a)}
\explique{On note $A(x\,;y)$ ; $I(1\,;1)$ est le milieu de $[AB]$ avec $B(5\,;-3)$.}{on
écrit la formule avec l'inconnue}
\begin{calculs}
\frac{x+5}{2} &= 1 &&\justif{abscisses}\\
x+5 &= 2 &&\justif{on multiplie par $2$}\\
x &= -3 &&
\end{calculs}
\begin{calculs}
\frac{y+(-3)}{2} &= 1 &&\justif{ordonnées}\\
y-3 &= 2 &&\\
y &= 5 &&
\end{calculs}
\explique{Donc $\res{A(-3\,;5)}$.}{contrôle : $\frac{-3+5}{2}=1$ et $\frac{5-3}{2}=1$}
\cas{b)}
\explique{Même calcul avec $B(-2\,;4)$ : $x-2=2$ et $y+4=2$, donc $\res{A(4\,;-2)}$.}{contrôle :
$\frac{4-2}{2}=1$ et $\frac{-2+4}{2}=1$}
\rem{Autre méthode : de $B$ à $I$, puis le même déplacement de $I$ à $A$.}
\end{exo}

\begin{exo}[De deux façons]
\cas{a)}
\begin{calculs}
\vect{EF} &\coord{1-(-2)}{2-5}=\coord{3}{-3} &&\justif{extrémité moins origine}\\
\vect{FG} &\coord{4-1}{-1-2}=\coord{3}{-3} &&
\end{calculs}
\explique{$\vect{EF}=\vect{FG}$, donc \res{$F$ est le milieu de $[EG]$}.}{chapitre 3,
théorème 3}
\cas{b)}
\explique{Milieu de $[EG]$ : $\left(\frac{-2+4}{2}\,;\frac{5+(-1)}{2}\right)$, soit
$\res{(1\,;2)}$ : c'est $F$.}{propriété 6}
\cas{c)}
\explique{Milieu de $[FH]$ : $\left(\frac{1+7}{2}\,;\frac{2+(-4)}{2}\right)$, soit
$(4\,;-1)$ : \res{oui}, c'est $G$.}{moyennes des coordonnées}
\cas{d)}
\begin{calculs}
\det(\vect{EF}\,;\vect{EH}) &= 3\times(-9)-(-3)\times 9 &&\justif{$\vect{EH}\coord{7-(-2)}{-4-5}=\coord{9}{-9}$}\\
\det(\vect{EF}\,;\vect{EH}) &= 0 &&\justif{$-27+27$}
\end{calculs}
\explique{$E$, $F$ et $H$ sont \res{alignés}.}{partie VI}
\end{exo}

\begin{exo}[La droite des milieux]
\cas{a)}
\explique{$I\left(\frac{-2+4}{2}\,;\frac{1+3}{2}\right)$, soit $\res{I(1\,;2)}$.}{milieu de
$[AB]$}
\explique{$J\left(\frac{-2+2}{2}\,;\frac{1+(-3)}{2}\right)$, soit $\res{J(0\,;-1)}$.}{milieu
de $[AC]$}
\cas{b)}
\explique{$\vect{IJ}\coord{-1}{-3}$ et $\vect{BC}\coord{-2}{-6}$.}{$\coord{0-1}{-1-2}$ et
$\coord{2-4}{-3-3}$}
\begin{calculs}
\det(\vect{IJ}\,;\vect{BC}) &= -1\times(-6)-(-3)\times(-2) &&\justif{$xy'-yx'$}\\
\det(\vect{IJ}\,;\vect{BC}) &= 0 &&\justif{$6-6$}
\end{calculs}
\explique{Les droites $(IJ)$ et $(BC)$ sont \res{parallèles}.}{vecteurs colinéaires}
\cas{c)}
\explique{Les coordonnées de $\vect{BC}$ sont le double de celles de $\vect{IJ}$ :
$\res{\vect{BC}=2\vect{IJ}}$.}{la droite des milieux est parallèle au troisième côté, et
deux fois plus courte (collège)}
\end{exo}

\partie{Norme d'un vecteur, longueur d'un segment}

\begin{exo}[Des normes]
\explique{$\norme{\vec{u}}=\sqrt{x^{2}+y^{2}}$ pour $\vec{u}\coord{x}{y}$, le repère étant
orthonormé.}{propriété 7}
\begin{calculs}
\norme{\vec{u}} &= \sqrt{0^{2}+(-1)^{2}}=\res{1} &&\justif{$\sqrt{1}$}\\
\norme{\vec{v}} &= \sqrt{3^{2}+4^{2}}=\res{5} &&\justif{$9+16=25$}\\
\norme{\vec{w}} &= \sqrt{8^{2}+(-6)^{2}}=\res{10} &&\justif{$64+36=100$}\\
\norme{\vec{t}} &= \sqrt{2^{2}+(-2)^{2}}=\sqrt{8} &&\justif{$4+4=8$}\\
\norme{\vec{t}} &= \res{2\sqrt{2}} &&\justif{$\sqrt{8}=\sqrt{4}\times\sqrt{2}$}\\
\norme{\vec{s}} &= \sqrt{(-1)^{2}+3^{2}}=\res{\sqrt{10}} &&\justif{$10$ n'a pas de facteur carré}
\end{calculs}
\rem{Parenthèses autour des nombres négatifs : $(-6)^{2}=36$, alors que $-6^{2}=-36$.}
\end{exo}

\begin{exo}[Des longueurs]
\begin{calculs}
AB &= \sqrt{3^{2}+4^{2}}=\res{5} &&\justif{$\vect{AB}\coord{4-1}{6-2}$ ; $9+16=25$}\\
CD &= \sqrt{6^{2}+(-8)^{2}}=\res{10} &&\justif{$\vect{CD}\coord{4-(-2)}{-5-3}$}\\
EF &= \sqrt{4^{2}+2^{2}}=\sqrt{20} &&\justif{$\vect{EF}\coord{4}{2}$}\\
EF &= \res{2\sqrt{5}} &&\justif{$\sqrt{20}=\sqrt{4}\times\sqrt{5}$}
\end{calculs}
\end{exo}

\begin{exo}[Nature d'un triangle]
\cas{a)}
\explique{$\vect{AB}\coord{4}{3}$, $\vect{BC}\coord{3}{-4}$ et $\vect{AC}\coord{7}{-1}$.}{extrémité
moins origine}
\begin{calculs}
AB^{2} &= 4^{2}+3^{2}=25 &&\justif{on calcule les carrés : pas de racine}\\
BC^{2} &= 3^{2}+(-4)^{2}=25 &&\\
AC^{2} &= 7^{2}+(-1)^{2}=50 &&\justif{le plus long côté}
\end{calculs}
\explique{$AB=BC$ et $AB^{2}+BC^{2}=50=AC^{2}$ : $ABC$ est \res{rectangle isocèle en
$B$}.}{isocèle ; réciproque du théorème de Pythagore}
\cas{b)}
\explique{$ABCD$ parallélogramme : $\vect{AB}=\vect{DC}$. Pour $D(x\,;y)$ : $6-x=4$ et
$1-y=3$, donc $\res{D(2\,;-2)}$.}{partie V}
\explique{Un parallélogramme qui a un angle droit (en $B$) et deux côtés consécutifs égaux
($AB=BC$) est \res{un carré}.}{collège}
\cas{c)}
\begin{calculs}
EF^{2} &= 4^{2}+1^{2}=17 &&\justif{$\vect{EF}\coord{4}{1}$}\\
EG^{2} &= 1^{2}+4^{2}=17 &&\justif{$\vect{EG}\coord{1}{4}$}\\
FG^{2} &= (-3)^{2}+3^{2}=18 &&\justif{$\vect{FG}\coord{-3}{3}$}
\end{calculs}
\explique{$EF=EG$ : $EFG$ est \res{isocèle en $E$}.}{deux côtés égaux}
\explique{$17+17=34\neq 18$ et $17+18=35\neq 17$ : \res{il n'est pas rectangle}.}{on teste
avec le plus grand côté, puis on vérifie les autres}
\end{exo}

\begin{exo}[Sur un cercle]
\explique{Un point est sur $\mathcal{C}$ quand sa distance à $K(1\,;-1)$ vaut $5$.}{on
calcule $KM$, $KN$, $KR$, $KS$}
\begin{calculs}
KM &= \sqrt{3^{2}+4^{2}}=5 &&\justif{$\vect{KM}\coord{3}{4}$}\\
KN &= \sqrt{(-3)^{2}+(-4)^{2}}=5 &&\justif{$\vect{KN}\coord{-3}{-4}$}\\
KR &= \sqrt{5^{2}+0^{2}}=5 &&\justif{$\vect{KR}\coord{5}{0}$}\\
KS &= \sqrt{(-1)^{2}+5^{2}}=\sqrt{26} &&\justif{$\vect{KS}\coord{-1}{5}$}
\end{calculs}
\explique{\res{$M$, $N$ et $R$ sont sur $\mathcal{C}$} ; \res{$S$ n'y est pas}.}{$\sqrt{26}\neq 5$}
\explique{Milieu de $[MN]$ : $\left(\frac{4+(-2)}{2}\,;\frac{3+(-5)}{2}\right)$, soit
$(1\,;-1)$ : c'est $K$. Donc \res{$[MN]$ est un diamètre} de $\mathcal{C}$.}{ses extrémités
sont sur le cercle et son milieu est le centre}
\end{exo}

\partie{Trigonométrie dans le triangle rectangle}

\begin{exo}[Calculer une longueur]
\cas{a)}
\explique{Hypoténuse $[AC]$ ; $[BC]$ est opposé à $\widehat{A}$, $[AB]$ lui est
adjacent.}{on repère les côtés avant de choisir le rapport}
\begin{calculs}
\sin 40^{\circ} &= \frac{BC}{10} &&\justif{opposé sur hypoténuse}\\
BC &= 10\times\sin 40^{\circ} &&\justif{produit en croix}\\
BC &\approx \res{6{,}4\text{~cm}} &&\justif{calculatrice en degrés}
\end{calculs}
\begin{calculs}
\cos 40^{\circ} &= \frac{AB}{10} &&\justif{adjacent sur hypoténuse}\\
AB &= 10\times\cos 40^{\circ} &&\\
AB &\approx \res{7{,}7\text{~cm}} &&
\end{calculs}
\cas{b)}
\explique{Hypoténuse $[DF]$ ; $[DE]$, connu, est adjacent à $\widehat{D}$ ; $[EF]$ lui est
opposé.}{opposé et adjacent : tangente (méthode 11 d)}
\begin{calculs}
\tan 25^{\circ} &= \frac{EF}{6} &&\justif{opposé sur adjacent}\\
EF &= 6\times\tan 25^{\circ} &&\justif{produit en croix}\\
EF &\approx \res{2{,}8\text{~cm}} &&
\end{calculs}
\begin{calculs}
\cos 25^{\circ} &= \frac{6}{DF} &&\justif{adjacent sur hypoténuse}\\
DF &= \frac{6}{\cos 25^{\circ}} &&\justif{l'inconnue est au dénominateur}\\
DF &\approx \res{6{,}6\text{~cm}} &&\justif{plus long que $DE$ : c'est l'hypoténuse}
\end{calculs}
\rem{Pour calculer $DF$, on utilise $DE$ (donné), pas $EF$ (arrondi) : on évite de
cumuler les arrondis.}
\end{exo}

\begin{exo}[Les côtés d'un triangle rectangle]
\cas{a)}
\explique{L'hypoténuse est \res{$[RT]$}, le côté opposé à l'angle droit. Pour l'angle
$\widehat{R}$, le côté adjacent est \res{$[RS]$} et le côté opposé est \res{$[ST]$}.}{le côté
adjacent touche le sommet $R$ ; l'opposé est en face}
\cas{b)}
\explique{Pour l'angle $\widehat{T}$, l'hypoténuse est toujours $[RT]$, le côté adjacent est
\res{$[ST]$} et le côté opposé est \res{$[RS]$}.}{adjacent et opposé s'échangent quand on
change d'angle}
\cas{c)}
\explique{$\res{\cos\widehat{R}=\dfrac{RS}{RT}}$, $\res{\sin\widehat{R}=\dfrac{ST}{RT}}$ et
$\res{\tan\widehat{T}=\dfrac{RS}{ST}}$.}{adjacent sur hypoténuse, opposé sur hypoténuse,
opposé sur adjacent}
\end{exo}

\begin{exo}[Calculer un angle]
\cas{a)}
\begin{calculs}
\cos\widehat{G} &= \frac{GH}{GK} &&\justif{$[GH]$ adjacent, $[GK]$ hypoténuse}\\
\cos\widehat{G} &= \frac{5}{8} &&\\
\widehat{G} &\approx \res{51^{\circ}} &&\justif{touche $\cos^{-1}$ : $\cos^{-1}(5\div 8)\approx 51{,}3$}
\end{calculs}
\cas{b)}
\begin{calculs}
\tan\widehat{L} &= \frac{MN}{LM} &&\justif{$[MN]$ opposé, $[LM]$ adjacent}\\
\tan\widehat{L} &= \frac{7}{4} &&\\
\widehat{L} &\approx \res{60^{\circ}} &&\justif{touche $\tan^{-1}$ : $\approx 60{,}3$}
\end{calculs}
\cas{c)}
\explique{$\vect{RT}\coord{4}{3}$, donc $RT=\sqrt{16+9}=\res{5}$.}{partie VIII}
\begin{calculs}
\cos\widehat{R} &= \frac{RS}{RT} &&\justif{$RS=4$ : $\vect{RS}\coord{4}{0}$}\\
\cos\widehat{R} &= \frac{4}{5} &&\\
\widehat{R} &\approx \res{37^{\circ}} &&\justif{$\cos^{-1}(0{,}8)\approx 36{,}9$}
\end{calculs}
\end{exo}

\begin{exo}[Une échelle]
\cas{a)}
\explique{L'échelle, le mur et le sol forment un triangle rectangle : l'échelle est
l'hypoténuse ($5$~m) ; la hauteur est opposée à l'angle de $70^{\circ}$, la distance au
mur lui est adjacente.}{le mur est vertical, le sol horizontal}
\begin{calculs}
h &= 5\times\sin 70^{\circ} &&\justif{$\sin 70^{\circ}=\frac{h}{5}$}\\
h &\approx \res{4{,}7\text{~m}} &&
\end{calculs}
\begin{calculs}
d &= 5\times\cos 70^{\circ} &&\justif{$\cos 70^{\circ}=\frac{d}{5}$}\\
d &\approx \res{1{,}7\text{~m}} &&
\end{calculs}
\cas{b)}
\begin{calculs}
\cos\alpha &= \frac{1{,}2}{5} &&\justif{adjacent sur hypoténuse}\\
\alpha &\approx 76^{\circ} &&\justif{$\cos^{-1}(0{,}24)\approx 76{,}1$}
\end{calculs}
\explique{$76^{\circ}>75^{\circ}$ : l'échelle est trop droite, elle \res{n'est pas posée en
sécurité}.}{on compare à l'intervalle demandé}
\end{exo}

\begin{exo}[Une rampe d'accès]
\cas{a)}
\begin{calculs}
\tan\alpha &= \frac{0{,}6}{8} &&\justif{opposé (la hauteur) sur adjacent (la longueur au sol)}\\
\alpha &\approx 4{,}3^{\circ} &&\justif{$\tan^{-1}(0{,}075)$}
\end{calculs}
\explique{$4{,}3^{\circ}>3^{\circ}$ : cette rampe \res{ne convient pas}, elle est trop
raide.}{on compare à l'angle maximal}
\cas{b)}
\begin{calculs}
\tan 3^{\circ} &= \frac{0{,}6}{L} &&\justif{angle maximal : $3^{\circ}$}\\
L &= \frac{0{,}6}{\tan 3^{\circ}} &&\justif{l'inconnue est au dénominateur}\\
L &\approx 11{,}45 &&
\end{calculs}
\explique{Il faut au moins \res{$11{,}5$~m} de longueur horizontale.}{on arrondit au-dessus :
avec $11{,}4$~m, l'angle dépasserait $3^{\circ}$}
\end{exo}

\end{document}
